\chapter{Filtro pasa altos}
\label{chap:pasa-altos}

\section{Especificaciones}
\label{sec:pa-especificaciones}

\begin{table}[!ht]
  \centering
  \begin{tabular}{|l|c|}
    \hline
    Parámetro & Valor \\
    \hline
    Orden & 4 \\
    Atenuación máxima en banda de paso $A_{max}$ & $\SI{2.25}{\decibel}$ \\
    Frecuencia de paso $f_p$ & $\SI{33}{\kilo\hertz}$ \\
    Ganancia en banda de paso & 12,5 V/V \\
    Aproximación & Butterworth \\
    Topología & Sallen-Key \\
    \hline
  \end{tabular}
  \caption{Especificaciones del filtro pasa altos.}
  \label{tab:pa-especificaciones}
\end{table}

\section{Diseño}
\label{sec:pa-diseno}

\subsection{Frecuencia de corte y polos de la aproximación}
\label{subsec:pa-polos}

El filtro pasa altos se obtiene del prototipo pasa bajos normalizado mediante la transformación
$s \rightarrow \omega_0 / s$, por lo que la frecuencia de corte a $\SI{-3}{\decibel}$ queda por debajo de $f_p$:

\begin{equation}
  f_c = f_p \left( 10^{A_{max}/10} - 1 \right)^{1/(2n)}
  \label{eq:pa-fc}
\end{equation}

% TODO: calcular f_c, listar los polos del Butterworth de orden 4 (Q_1 = 0,5412 y Q_2 = 1,3066) y repartir la ganancia
% total de 12,5 V/V entre las dos etapas. Verificar el producto ganancia-ancho de banda del operacional elegido
% (ganancia por etapa x frecuencia máxima de interés).

\subsection{Cálculo de las etapas Sallen-Key}
\label{subsec:pa-etapas}

Cada etapa de segundo orden (figura~\ref{crkt:sk-pasa-altos}) tiene ganancia $K = 1 + R_b / R_a$ y

\begin{align*}
  \omega_0 &= \frac{1}{\sqrt{R_1 R_2 C_1 C_2}} \\
  Q &= \frac{\sqrt{R_1 R_2 C_1 C_2}}{R_2 C_1 + R_2 C_2 + R_1 C_2 \left( 1 - K \right)}
\end{align*}

\begin{figure}[!ht]
  \centering
  \begin{tikzpicture}[american]
  \node[op amp, noinv input up] (oa) at (6,0) {};

  \draw (0,0.5) node[ocirc, label=left:$V_i$] {} to[C, l=$C_1$] (1.8,0.5) node[circ] (n1) {}
    to[C, l=$C_2$] (3.6,0.5) node[circ] (n2) {} -- (oa.+);
  \draw (n2) to[R, l_=$R_2$] ++(0,-2) node[ground] {};
  \draw (oa.out) -- ++(0.6,0) node[circ] (out) {} to[short, -o] ++(1,0) node[right] {$V_o$};
  \draw (n1) to[R, l=$R_1$] ++(0,2) -| (out);

  % Red de ganancia K = 1 + Rb/Ra
  \draw (oa.-) -- ++(-0.5,0) coordinate (m) -- ++(0,-1.1) node[circ] (m2) {}
    to[R, l=$R_a$] ++(0,-1.8) node[ground] {};
  \draw (m2) to[R, l_=$R_b$] (out |- m2) -- (out);
\end{tikzpicture}

  \caption{Etapa Sallen-Key pasa altos de segundo orden con ganancia $K$.}
  \label{crkt:sk-pasa-altos}
\end{figure}

% TODO: desarrollar el cálculo de R_1, R_2, C_1, C_2, R_a y R_b de cada etapa.

\subsection{Valores comerciales}
\label{subsec:pa-comerciales}

\begin{table}[!ht]
  \centering
  \begin{tabular}{|l|c|c|c|c|}
    \hline
     & \multicolumn{2}{c|}{Etapa 1} & \multicolumn{2}{c|}{Etapa 2} \\
    \cline{2-5}
    Componente & Calculado & Comercial & Calculado & Comercial \\
    \hline
    $R_1$ & & & & \\
    $R_2$ & & & & \\
    $C_1$ & & & & \\
    $C_2$ & & & & \\
    $R_a$ & & & & \\
    $R_b$ & & & & \\
    \hline
  \end{tabular}
  \caption{Valores calculados y comerciales del filtro pasa altos.}
  \label{tab:pa-componentes}
\end{table}

\section{Simulación}
\label{sec:pa-simulacion}

\subsection{Función de transferencia con elementos ideales}
\label{subsec:pa-sim-ideal}

% TODO: diagrama de Bode (módulo y fase) con los valores calculados.

\subsection{Función de transferencia con valores comerciales}
\label{subsec:pa-sim-real}

% TODO: repetir con los valores comerciales y con el modelo real del operacional; comparar con las especificaciones.

\section{Implementación y mediciones}
\label{sec:pa-mediciones}

% TODO: si este filtro es uno de los dos implementados: armado, fotos en pictures/, tabla de barrido en frecuencia,
% captura del perfil de frecuencia (FFT) del osciloscopio y comparación medido vs. simulado. Si no se implementa,
% indicar que se realiza en el TP de filtros digitales.

\begin{table}[!ht]
  \centering
  \begin{tabular}{|c|c|c|c|c|}
    \hline
    $f$ [kHz] & $V_{in,pp}$ [V] & $V_{out,pp}$ [V] & $|H|$ medido [dB] & $|H|$ simulado [dB] \\
    \hline
     & & & & \\
     & & & & \\
     & & & & \\
    \hline
  \end{tabular}
  \caption{Relevamiento de la respuesta en frecuencia del filtro pasa altos.}
  \label{tab:pa-barrido}
\end{table}
